\answer{ex:17:1}{$0.2$~W contre $81$~µW; le flux brut exigerait au plus $50$~pJ/bit.}
\answer{ex:17:2}{Pour $c=0.1$: $5.6$, $8.7$, $9.2$~bit/s, limite $9.3$~bit/s; pour $c=0$ et $N=1000$: $65$~bit/s.}
\answer{ex:17:3}{$B=\log_2N+P\log_2P+(1-P)\log_2\frac{1-P}{N-1}=4.87$, $4.37$, $3.29$~bits par caractère, soit $7.3$, $6.6$, $4.9$~bit/s; pour $10$~bit/s, il faut $2.3$~caractères par seconde.}
\answer{ex:17:4}{Variance de l'erreur $2b/(Nm^2T)+\sigma_\xi^2$ par composante; les contributions sont égales pour $N\approx90$; limite $\tfrac12\log_2(1+0.25/0.04)\approx1.4$~bit par composante et par fenêtre.}
\answer{ex:17:5}{La paire converge si $\eta_bd^2+\eta_db^2<2$, soit à peu près pour $\eta<1$: à $0.8$ elle converge, à $1.1$ oscillations entretenues de période $2$, à $1.5$ oscillations non périodiques, à $2$ divergence; il faut séparer les vitesses des élèves.}
\answer{ex:17:6}{Seuil $\approx4.4\sigma$; $102$~kbit/s, $0.10$~mW contre $205$~mW pour le flux brut, soit $2000$ fois moins; seuls $5.7\cdot10^{-5}$ des échantillons saturent (plus de $3$ impulsions).}
\answer{ex:17:7}{Question ouverte. Environ $46$~bit/s pour $\text{SNR}\approx0.9$ et environ $330$~bit/s avec un bruit indépendant. Si c'est un bruit semblable au signal qui gêne, l'information extraite sature quand le nombre de canaux augmente; si c'est la méconnaissance du code, elle augmente avec un meilleur décodeur (par exemple, qui tient compte des instants des impulsions).}
